\documentclass[11pt]{article}
\usepackage[margin=2.5cm]{geometry}
\usepackage{amsmath, amssymb, physics}
\usepackage{hyperref}

\newcommand{\bk}{\mathbf{k}}
\newcommand{\bq}{\mathbf{q}}
\newcommand{\br}{\mathbf{r}}
\newcommand{\GR}{G^{R}}
\newcommand{\GA}{G^{A}}

\title{LDOS modulation $\delta\rho$ from the Born $T$-matrix sum $S(\bq)$}
\date{}

\begin{document}
\maketitle

These notes follow the conventions of \texttt{QPI.pdf}: $S(\bq)$ is built from $g_0(\bk)\,T\,g_0(\bk-\bq)$, and the FT-LDOS is $\rho(\bq) = \sum_{\br} e^{-i\bq\cdot\br}\rho(\br)$. On the lattice, integrals $\int d^2k/(2\pi)^2$ become $\frac{1}{N}\sum_{\bk}$, and the impurity potential is allowed to depend on momentum, $T \to T(\bk,\bk')$.

\section{Conventions}

The lattice has $N$ sites $\br$ and $N$ momenta $\bk$ in the Brillouin zone. Operators are related by
\begin{equation}
  \Psi_{\br} = \frac{1}{\sqrt{N}} \sum_{\bk} e^{i\bk\cdot\br}\, \Psi_{\bk},
  \qquad
  \Psi_{\bk} = \begin{pmatrix} c_{\bk\uparrow} \\ c^{\dagger}_{-\bk\downarrow} \end{pmatrix}.
\end{equation}
Both Nambu components of $\Psi_{\br}$ carry the same phase, since $c^{\dagger}_{\br\downarrow} = N^{-1/2}\sum_{\bk} e^{i\bk\cdot\br} c^{\dagger}_{-\bk\downarrow}$.

The Green's function is built from $\Psi_{\br}$ and $\Psi^{\dagger}_{\br'}$, e.g.\ $\GR(\br,\br';t) = -i\theta(t)\langle\{\Psi_{\br}(t), \Psi^{\dagger}_{\br'}(0)\}\rangle$. The creation operator is the Hermitian conjugate of the same transform, $\Psi^{\dagger}_{\br'} = N^{-1/2}\sum_{\bk'} e^{-i\bk'\cdot\br'}\Psi^{\dagger}_{\bk'}$, so the two ends carry opposite phases:
\begin{equation}
  G(\br,\br') = \frac{1}{N}\sum_{\bk,\bk'} e^{i\bk\cdot\br}\, G(\bk,\bk')\, e^{-i\bk'\cdot\br'} .
  \label{eq:G_rr}
\end{equation}
This is the change of basis $\langle\br|G|\br'\rangle = \sum \langle\br|\bk\rangle\langle\bk|G|\bk'\rangle\langle\bk'|\br'\rangle$ with $\langle\br|\bk\rangle = e^{i\bk\cdot\br}/\sqrt{N}$. As a check, a translation-invariant $G(\bk,\bk') = \delta_{\bk\bk'}\,g_0(\bk)$ gives $G_0(\br,\br') = \frac{1}{N}\sum_{\bk} e^{i\bk\cdot(\br-\br')}\,g_0(\bk)$, which depends only on $\br-\br'$ (cf.\ eq.~(13) of \texttt{QPI.pdf}).

The bare retarded and advanced Green's functions are
\begin{equation}
  g^R_0(\bk,\omega) = \left[\omega + i\gamma - H_{\mathrm{BdG}}(\bk)\right]^{-1},
  \qquad
  g^A_0(\bk,\omega) = \left[g^R_0(\bk,\omega)\right]^{\dagger},
\end{equation}
and in general $\GA(\br,\br') = [\GR(\br',\br)]^{\dagger}$, so on the diagonal $\GA(\br,\br) = [\GR(\br,\br)]^{\dagger}$.

\section{Born approximation and $S(\bq)$}

For an impurity $H_{\mathrm{imp}} = \sum_{\bk\bk'} \Psi^{\dagger}_{\bk} T(\bk,\bk') \Psi_{\bk'}$ with $T(\bk,\bk')^{\dagger} = T(\bk',\bk)$, the first-order correction is
\begin{equation}
  \delta\GR(\bk,\bk') = g^R_0(\bk)\, T(\bk,\bk')\, g^R_0(\bk') .
\end{equation}
Following eq.~(21) of \texttt{QPI.pdf}, define
\begin{equation}
  S(\bq) = \sum_{\bk} g^R_0(\bk)\, T(\bk,\bk-\bq)\, g^R_0(\bk-\bq) = \sum_{\bk} \delta\GR(\bk,\bk-\bq),
  \label{eq:S}
\end{equation}
a $2\times2$ Nambu matrix. This is what the code computes (\texttt{g0k}, \texttt{T}, \texttt{g0kmq}). Both Green's functions in $S$ are retarded.

\section{Real-space diagonal}

Setting $\br' = \br$ in \eqref{eq:G_rr} and writing $\bk' = \bk - \bq$:
\begin{equation}
  \delta\GR(\br,\br) = \frac{1}{N}\sum_{\bk,\bq} e^{i\bk\cdot\br}\, \delta\GR(\bk,\bk-\bq)\, e^{-i(\bk-\bq)\cdot\br}
  = \frac{1}{N}\sum_{\bq} e^{+i\bq\cdot\br}\, S(\bq),
  \label{eq:dGR}
\end{equation}
so $S(\bq) = \sum_{\br} e^{-i\bq\cdot\br}\,\delta\GR(\br,\br)$ is the standard Fourier transform of $\delta\GR(\br,\br) = S(\br)$, as in eqs.~(28)--(29) of \texttt{QPI.pdf}.

The advanced part is the Hermitian conjugate of \eqref{eq:dGR}:
\begin{equation}
  \delta\GA(\br,\br) = \left[\delta\GR(\br,\br)\right]^{\dagger}
  = \frac{1}{N}\sum_{\bq} e^{-i\bq\cdot\br}\, S(\bq)^{\dagger}
  = \frac{1}{N}\sum_{\bq} e^{+i\bq\cdot\br}\, S(-\bq)^{\dagger}.
  \label{eq:dGA}
\end{equation}
$S(-\bq)^{\dagger}$ is the advanced counterpart of $S(\bq)$ with the same momentum structure. Using $T(\bk,\bk')^{\dagger} = T(\bk',\bk)$ and shifting $\bk \to \bk-\bq$ (cf.\ eqs.~(22)--(23) of \texttt{QPI.pdf}),
\begin{equation}
  S(-\bq)^{\dagger} = \sum_{\bk} g^A_0(\bk+\bq)\, T(\bk+\bq,\bk)\, g^A_0(\bk)
  = \sum_{\bk} g^A_0(\bk)\, T(\bk,\bk-\bq)\, g^A_0(\bk-\bq)
  = \sum_{\bk} \delta\GA(\bk,\bk-\bq).
\end{equation}

\section{FT-LDOS}

The spin-up electron LDOS is
\begin{equation}
  \delta\rho(\br,\omega) = -\frac{1}{\pi}\,\Im\, \delta\GR_{00}(\br,\br,\omega)
  = -\frac{1}{2\pi i}\left[\delta\GR(\br,\br) - \delta\GA(\br,\br)\right]_{00}.
\end{equation}
Substituting \eqref{eq:dGR} and \eqref{eq:dGA} gives $\delta\rho(\br) = \frac{1}{N}\sum_{\bq} e^{i\bq\cdot\br}\,\delta\rho(\bq)$ with
\begin{equation}
  \boxed{\;
  \delta\rho(\bq) = \sum_{\br} e^{-i\bq\cdot\br}\,\delta\rho(\br)
  = -\frac{1}{2\pi i}\left[S(\bq) - S(-\bq)^{\dagger}\right]_{00} , \;}
  \label{eq:rho}
\end{equation}
which is eq.~(24) of \texttt{QPI.pdf}. The inverse transform $\delta\rho(\br) = \frac{1}{N_q^2}\sum_{\bq} e^{+i\bq\cdot\br}\,\delta\rho(\bq)$ is \texttt{numpy.fft.ifft2}, as used in the code.

\paragraph{Relation to the earlier $\bk+\bq$ form.}
An earlier version of the code used $S_{+}(\bq) = \sum_{\bk}\delta\GR(\bk,\bk+\bq) = S(-\bq)$. In terms of $S_{+}$, \eqref{eq:rho} reads $\delta\rho(\bq) = \frac{i}{2\pi}\left[S_{+}(-\bq) - S_{+}(\bq)^{\dagger}\right]_{00}$. That is the same expression with $\bq$ relabelled, and it is the source of the apparent sign difference in $\bq$. The two codes give bitwise-identical $\delta\rho$.

\paragraph{Retarded vs.\ advanced.}
Both terms in \eqref{eq:rho} contain one retarded and one advanced function at the same $\bq$. Swapping $\GR \leftrightarrow \GA$ throughout flips the overall sign of $\delta\rho$, since $-\frac{1}{\pi}\Im\GA = +\frac{1}{\pi}\Im\GR$. It does not change the sign of $\bq$.

\section{Consequences for the symmetry analysis}

Because $\delta\rho(\br)$ is real, $\delta\rho(-\bq) = \delta\rho(\bq)^{*}$; this is manifest in \eqref{eq:rho}. Hence $\Re\,\delta\rho(\bq)$ is even and $\Im\,\delta\rho(\bq)$ is odd under $\bq\to-\bq$. In two dimensions inversion is the rotation $C_2 \in C_{6v}$, with $\chi(C_2) = +d_\nu$ for $A_1, A_2, E_2$ and $-d_\nu$ for $B_1, B_2, E_1$. So
\begin{equation}
  \Re\,\delta\rho(\bq) \in A_1 \oplus A_2 \oplus E_2,
  \qquad
  \Im\,\delta\rho(\bq) \in B_1 \oplus B_2 \oplus E_1,
  \qquad
  \Im\,\delta\rho(\br) = 0 .
\end{equation}

\section{Implementation notes}

\begin{itemize}
  \item On the centred $q$-mesh, $\bq = (m_1\mathbf{b}_1 + m_2\mathbf{b}_2)/N_q$ sits at index $(m_2 + N_q/2,\, m_1 + N_q/2)$. Mesh axis 0 steps along $\mathbf{b}_2$ and axis 1 along $\mathbf{b}_1$. $g_0(\bk-\bq)$ is \texttt{np.roll} of $g_0(\bk)$ by $+(m_2, m_1)\,N_k/N_q$.
  \item $V(\bk,\bk-\bq)$ is evaluated with $\bk$ and $\bk-\bq$ each folded into the Wigner--Seitz cell, which keeps $V$ Hermitian and $C_{6v}$-covariant.
  \item $S(-\bq)$ does not need a separate calculation. $-\bq$ sits at index $(N_q - i) \bmod N_q$ along each axis. Since $g_0$ is periodic in $\bk$ and $V$ uses folded momenta, $S(\bq + \mathbf{G}) = S(\bq)$, so the reindexed array is exact. This matches a direct recalculation of $S(-\bq)$ to $\sim 10^{-16}$ relative.
  \item The overall factor $1/N$ is omitted in the code (it only sets the scale).
\end{itemize}

\section{Numerical check}

On a $12\times12$ lattice ($N_q = N_k = 12$, $d+id$ pairing, $\omega = 1.5\,\Delta_{\mathrm{gap}}$, $\gamma = 10^{-2}$), with a generic Hermitian $V(\bk,\bk')$ that has no $C_{6v}$ symmetry:
\begin{itemize}
  \item $\delta\GR_{00}(\bk,\bk')$ was built in the full $(\bk,\bk')$ basis, rotated to sites using $\langle\br|\bk\rangle = e^{i\bk\cdot\br}/\sqrt{N}$, and $-\frac{1}{\pi}\Im\,\delta\GR_{00}(\br,\br)$ taken directly.
  \item The code's path (\eqref{eq:S}, \eqref{eq:rho}, then \texttt{ifft2}) reproduces it with ratio $1 \pm 10^{-12}$, and its imaginary part is $\sim 10^{-16}$.
  \item $S(\bq)$ equals $\sum_{\br} e^{-i\bq\cdot\br}\,\delta\GR_{00}(\br,\br)$ to $2\times10^{-15}$ relative.
  \item Replacing $g^R_0 \to g^A_0$ in $S$ gives exactly $-\delta\rho(\br)$.
\end{itemize}

\end{document}
